% Options for packages loaded elsewhere
% Options for packages loaded elsewhere
\PassOptionsToPackage{unicode}{hyperref}
\PassOptionsToPackage{hyphens}{url}
\PassOptionsToPackage{dvipsnames,svgnames,x11names}{xcolor}
%
\documentclass[
  letterpaper,
  DIV=11,
  numbers=noendperiod]{scrartcl}
\usepackage{xcolor}
\usepackage{amsmath,amssymb}
\setcounter{secnumdepth}{-\maxdimen} % remove section numbering
\usepackage{iftex}
\ifPDFTeX
  \usepackage[T1]{fontenc}
  \usepackage[utf8]{inputenc}
  \usepackage{textcomp} % provide euro and other symbols
\else % if luatex or xetex
  \usepackage{unicode-math} % this also loads fontspec
  \defaultfontfeatures{Scale=MatchLowercase}
  \defaultfontfeatures[\rmfamily]{Ligatures=TeX,Scale=1}
\fi
\usepackage{lmodern}
\ifPDFTeX\else
  % xetex/luatex font selection
\fi
% Use upquote if available, for straight quotes in verbatim environments
\IfFileExists{upquote.sty}{\usepackage{upquote}}{}
\IfFileExists{microtype.sty}{% use microtype if available
  \usepackage[]{microtype}
  \UseMicrotypeSet[protrusion]{basicmath} % disable protrusion for tt fonts
}{}
\makeatletter
\@ifundefined{KOMAClassName}{% if non-KOMA class
  \IfFileExists{parskip.sty}{%
    \usepackage{parskip}
  }{% else
    \setlength{\parindent}{0pt}
    \setlength{\parskip}{6pt plus 2pt minus 1pt}}
}{% if KOMA class
  \KOMAoptions{parskip=half}}
\makeatother
% Make \paragraph and \subparagraph free-standing
\makeatletter
\ifx\paragraph\undefined\else
  \let\oldparagraph\paragraph
  \renewcommand{\paragraph}{
    \@ifstar
      \xxxParagraphStar
      \xxxParagraphNoStar
  }
  \newcommand{\xxxParagraphStar}[1]{\oldparagraph*{#1}\mbox{}}
  \newcommand{\xxxParagraphNoStar}[1]{\oldparagraph{#1}\mbox{}}
\fi
\ifx\subparagraph\undefined\else
  \let\oldsubparagraph\subparagraph
  \renewcommand{\subparagraph}{
    \@ifstar
      \xxxSubParagraphStar
      \xxxSubParagraphNoStar
  }
  \newcommand{\xxxSubParagraphStar}[1]{\oldsubparagraph*{#1}\mbox{}}
  \newcommand{\xxxSubParagraphNoStar}[1]{\oldsubparagraph{#1}\mbox{}}
\fi
\makeatother


\usepackage{longtable,booktabs,array}
\usepackage{calc} % for calculating minipage widths
% Correct order of tables after \paragraph or \subparagraph
\usepackage{etoolbox}
\makeatletter
\patchcmd\longtable{\par}{\if@noskipsec\mbox{}\fi\par}{}{}
\makeatother
% Allow footnotes in longtable head/foot
\IfFileExists{footnotehyper.sty}{\usepackage{footnotehyper}}{\usepackage{footnote}}
\makesavenoteenv{longtable}
\usepackage{graphicx}
\makeatletter
\newsavebox\pandoc@box
\newcommand*\pandocbounded[1]{% scales image to fit in text height/width
  \sbox\pandoc@box{#1}%
  \Gscale@div\@tempa{\textheight}{\dimexpr\ht\pandoc@box+\dp\pandoc@box\relax}%
  \Gscale@div\@tempb{\linewidth}{\wd\pandoc@box}%
  \ifdim\@tempb\p@<\@tempa\p@\let\@tempa\@tempb\fi% select the smaller of both
  \ifdim\@tempa\p@<\p@\scalebox{\@tempa}{\usebox\pandoc@box}%
  \else\usebox{\pandoc@box}%
  \fi%
}
% Set default figure placement to htbp
\def\fps@figure{htbp}
\makeatother





\setlength{\emergencystretch}{3em} % prevent overfull lines

\providecommand{\tightlist}{%
  \setlength{\itemsep}{0pt}\setlength{\parskip}{0pt}}



 


\KOMAoption{captions}{tableheading}
\makeatletter
\@ifpackageloaded{caption}{}{\usepackage{caption}}
\AtBeginDocument{%
\ifdefined\contentsname
  \renewcommand*\contentsname{Table of contents}
\else
  \newcommand\contentsname{Table of contents}
\fi
\ifdefined\listfigurename
  \renewcommand*\listfigurename{List of Figures}
\else
  \newcommand\listfigurename{List of Figures}
\fi
\ifdefined\listtablename
  \renewcommand*\listtablename{List of Tables}
\else
  \newcommand\listtablename{List of Tables}
\fi
\ifdefined\figurename
  \renewcommand*\figurename{Figure}
\else
  \newcommand\figurename{Figure}
\fi
\ifdefined\tablename
  \renewcommand*\tablename{Table}
\else
  \newcommand\tablename{Table}
\fi
}
\@ifpackageloaded{float}{}{\usepackage{float}}
\floatstyle{ruled}
\@ifundefined{c@chapter}{\newfloat{codelisting}{h}{lop}}{\newfloat{codelisting}{h}{lop}[chapter]}
\floatname{codelisting}{Listing}
\newcommand*\listoflistings{\listof{codelisting}{List of Listings}}
\makeatother
\makeatletter
\makeatother
\makeatletter
\@ifpackageloaded{caption}{}{\usepackage{caption}}
\@ifpackageloaded{subcaption}{}{\usepackage{subcaption}}
\makeatother
\makeatletter
\@ifpackageloaded{tcolorbox}{}{\usepackage{tcolorbox}}
\makeatother
\makeatletter
\@ifpackageloaded{fontawesome5}{}{\usepackage{fontawesome5}}
\makeatother
\usepackage{bookmark}
\IfFileExists{xurl.sty}{\usepackage{xurl}}{} % add URL line breaks if available
\urlstyle{same}
\hypersetup{
  pdftitle={Quarto Extension},
  pdfauthor={Muhammad N. Ahmad},
  colorlinks=true,
  linkcolor={blue},
  filecolor={Maroon},
  citecolor={Blue},
  urlcolor={Blue},
  pdfcreator={LaTeX via pandoc}}


\title{Quarto Extension}
\author{Muhammad N. Ahmad}
\date{2025-10-31}
\begin{document}
\maketitle


\begin{centering}\begin{tcolorbox}[hbox,
colframe=lightgray,
colback=white]

\textcolor{Orchid}{{\Large {\faMastodon}}} %

\href{https://sciencemastodon.com/@gdyson/111371030805595929}{\raisebox{0.1 em}{Click to view embedded Mastodon post.}}

\end{tcolorbox}\end{centering}
\vspace{1em}

Linkedin \begin{centering}\begin{tcolorbox}[hbox,
colframe=lightgray,
colback=white]
\\
\textcolor{RoyalBlue}{{\Large {\faLinkedin}}} %\\
\href{https://www.linkedin.com/feed/update/urn:li:activity:7335234505370169344}{\raisebox{0.1 em}{Click to view embedded LinkedIn post.}}\\

\end{tcolorbox}\end{centering}
\vspace{1em}

Twitter

Twitter's option

For UKPack NbS project (funded by UKPack), @CIFOR\_ICRAF conducted a
capacity building workshop at county level to leverage citizen science
for critical data collection to help in decision
making.@OfficialMakueni, @UKinKenya, @lawinowiecki\#RegreeningApp
pic.twitter.com/75BGH6BCBb

--- Muhammad Ahmad (@mnabiahmad) May 15, 2024



<div class= "page-columns page-rows-contents page-layout-article"><div class="social-share"><a href="https://twitter.com/share?url=https://mnahmad.github.io/scriptndebug/posts/regex/index.html&text=Regex 101" target="_blank" class="twitter"><i class="fab fa-twitter fa-fw fa-lg"></i></a><a href="https://www.linkedin.com/shareArticle?url=https://mnahmad.github.io/scriptndebug/posts/regex/index.html&title=Regex 101" target="_blank" class="linkedin"><i class="fa-brands fa-linkedin-in fa-fw fa-lg"></i></a>  <a href="mailto:?subject=Regex 101&body=Check out this link:https://mnahmad.github.io/scriptndebug/posts/regex/index.html" target="_blank" class="email"><i class="fa-solid fa-envelope fa-fw fa-lg"></i></a><a href="https://www.facebook.com/sharer.php?u=https://mnahmad.github.io/scriptndebug/posts/regex/index.html" target="_blank" class="facebook"><i class="fab fa-facebook-f fa-fw fa-lg"></i></a><a href="https://reddit.com/submit?url=https://mnahmad.github.io/scriptndebug/posts/regex/index.html&title=Regex 101" target="_blank" class="reddit">   <i class="fa-brands fa-reddit-alien fa-fw fa-lg"></i></a><a href="https://www.stumbleupon.com/submit?url=https://mnahmad.github.io/scriptndebug/posts/regex/index.html&title=Regex 101" target="_blank" class="stumbleupon"><i class="fa-brands fa-stumbleupon fa-fw fa-lg"></i></a><a href="https://www.tumblr.com/share/link?url=https://mnahmad.github.io/scriptndebug/posts/regex/index.html&name=Regex 101" target="_blank" class="tumblr"><i class="fa-brands fa-tumblr fa-fw fa-lg"></i></a><a href="javascript:void(0);" onclick="var mastodon_instance=prompt('Mastodon Instance / Server Name?'); if(typeof mastodon_instance==='string' &amp;&amp; mastodon_instance.length){this.href='https://'+mastodon_instance+'/share?text=Regex 101 https://mnahmad.github.io/scriptndebug/posts/regex/index.html'}else{return false;}" target="_blank" class="mastodon"><i class="fa-brands fa-mastodon fa-fw fa-lg"></i></a><a href="https://bsky.app/intent/compose?text=https://mnahmad.github.io/scriptndebug/posts/regex/index.html Regex 101" target="_blank" class="bsky"><i class="fa-brands fa-bluesky"></i></a></div></div>


\end{document}
